\section{Conserving BCS density response}\label{app:D.21}

The independent Bogoliubov quasiparticles of \cref{app:D.20} respond to a uniform potential: at $\vect{q} = 0$ the density creates pairs of them with weight $\Delta_{\mathrm{BCS}}^2/E_{\vect{k}}^2$, although the electron number is conserved and a uniform potential can excite nothing. The defect is not in the mean-field ground state but in the response computed from it. When the density perturbs the quasiparticles, the pair field $\Delta_{\mathrm{BCS}}$ must respond as well, and its phase fluctuations restore charge conservation. This section writes the problem in Nambu form~\cite{nambu1960quasi} and shows that the bare bubble violates the continuity equation by a term proportional to $\Delta_{\mathrm{BCS}}$. The conserving response of Anderson~\cite{anderson1958random} is then derived as the linear response of the time-dependent mean field. We prove that it satisfies the continuity equation and the f-sum rule, and state the Ward identity that its vertex obeys. The result is used in \cref{sec:7.2}.

\subsection*{Nambu form}

Collect the two operators that the pairing mixes into the spinor $\Psi_{\vect{k}} = (c_{\vect{k}\uparrow}, c^\dagger_{-\vect{k}\downarrow})^{\mathrm{T}}$, and use the matrices
\begin{equation*}
  \boldsymbol{\eta} = \begin{pmatrix}1 & 0\\ 0 & -1\end{pmatrix}, \qquad
  \mathsf{s}_1 = \begin{pmatrix}0 & 1\\ 1 & 0\end{pmatrix}, \qquad
  \mathsf{s}_2 = \begin{pmatrix}0 & -\ii\\ \ii & 0\end{pmatrix} ,
\end{equation*}
with $\boldsymbol{\eta}$ the matrix of \cref{app:C.7}; together with the unit matrix they obey the Pauli algebra, $\mathsf{s}_1\mathsf{s}_2 = \ii\boldsymbol{\eta}$ and cyclically. For a $2\times2$ matrix $A$, define $O_A(\vect{q}) = \sum_{\vect{k}}\Psi^\dagger_{\vect{k}+\vect{q}}A\,\Psi_{\vect{k}}$. Writing out the components, $O_{\boldsymbol{\eta}}(\vect{q}) = \sum_{\vect{k}}[c^\dagger_{\vect{k}+\vect{q}\uparrow}c_{\vect{k}\uparrow} - c_{-\vect{k}-\vect{q}\downarrow}c^\dagger_{-\vect{k}\downarrow}] = n_e(\vect{q})$ for $\vect{q} \ne 0$, after anticommuting and relabeling $-\vect{k} - \vect{q} \to \vect{k}$ in the spin-down term. The density is the $\boldsymbol{\eta}$ vertex. The pair operators are $O_{\mathsf{s}_1}(\vect{q}) = B^\dagger_{\vect{q}} + B_{-\vect{q}}$ and $O_{\mathsf{s}_2}(\vect{q}) = -\ii(B^\dagger_{\vect{q}} - B_{-\vect{q}})$, with $B^\dagger_{\vect{P}} = \sum_{\vect{k}}c^\dagger_{\vect{k}+\vect{P}\uparrow}c^\dagger_{-\vect{k}\downarrow}$ the creation operator of a pair of total momentum $\vect{P}$. For a real gap, $\mathsf{s}_1$ fluctuations change the magnitude of the pair field and $\mathsf{s}_2$ fluctuations its phase.

The mean-field Hamiltonian of \cref{app:D.20} is $H_{\mathrm{MF}} = \sum_{\vect{k}}\Psi^\dagger_{\vect{k}}h_{\vect{k}}\Psi_{\vect{k}} + \text{const}$ with $h_{\vect{k}} = \xi_{\vect{k}}\boldsymbol{\eta} - \Delta_{\mathrm{BCS}}\mathsf{s}_1$, whose eigenvalues are $\pm E_{\vect{k}}$, with projectors $P^\pm_{\vect{k}} = \frac12(1 \pm h_{\vect{k}}/E_{\vect{k}})$. The Matsubara Green's function is the $2\times2$ matrix $\GF_0(\vect{k},\ii\omega^{\mathrm{F}}) = (\ii\omega^{\mathrm{F}} - h_{\vect{k}})^{-1} = \sum_\pm P^\pm_{\vect{k}}/(\ii\omega^{\mathrm{F}} \mp E_{\vect{k}})$. The off-diagonal elements are the anomalous averages of the pair state.

\subsection*{The interaction in the pairing channel}

To let the pair field respond at wavevector $\vect{q}$, the interaction must scatter pairs of every total momentum, not only $\vect{P} = 0$ as in the reduced Hamiltonian \cref{eq:D.20-H}. Take the contact attraction
\begin{equation*}
  H_{\mathrm{int}} = -\frac{V_{\mathrm{BCS}}}{\Vol}\sum_{\vect{P}}B^\dagger_{\vect{P}}B_{\vect{P}} = -V_{\mathrm{BCS}}\int\dd^3r\,\psi^\dagger_\uparrow\psi^\dagger_\downarrow\psi_\downarrow\psi_\uparrow(\vect{r}) ,
\end{equation*}
with the pair states restricted to the shell $|\xi| < \omega_{\mathrm{D}}$. Without the restriction it is a function of the local densities $\psi^\dagger_\sigma\psi_\sigma$ alone and commutes with the density. On a lattice it is the on-site term $-V_{\mathrm{BCS}}\sum_in_{i\uparrow}n_{i\downarrow}$, which needs no shell; in the continuum the shell makes it slightly nonlocal, and its commutator with $n_e(\vect{q})$ is of relative order $v_{\mathrm{F}}q/\omega_{\mathrm{D}}$, negligible at the wavevectors of interest. The Hamiltonian $H_0 + H_{\mathrm{int}}$ therefore obeys the continuity equation \cref{eq:4.3-continuity}, $[H, n_e(\vect{q})] = \vect{q}\cdot\vect{j}_e(\vect{q})$. With $B^\dagger_{\vect{P}} = \frac12(O_{\mathsf{s}_1} + \ii O_{\mathsf{s}_2})(\vect{P})$ and $B_{\vect{P}} = \frac12(O_{\mathsf{s}_1} - \ii O_{\mathsf{s}_2})(-\vect{P})$,
\begin{equation*}
  H_{\mathrm{int}} = -\frac{V_{\mathrm{BCS}}}{4\Vol}\sum_{\vect{P}}\sum_{a=1,2}O_{\mathsf{s}_a}(\vect{P})\,O_{\mathsf{s}_a}(-\vect{P}) + (\text{one-body terms}) ,
\end{equation*}
where the one-body terms come from reordering and shift the chemical potential. Replacing one factor of each product by its average, in the two possible ways, gives the mean-field term $-(V_{\mathrm{BCS}}/2\Vol)\sum_a\langle O_{\mathsf{s}_a}(\vect{0})\rangle O_{\mathsf{s}_a}(\vect{0})$, and with $\langle O_{\mathsf{s}_1}(\vect{0})\rangle = 2\sum_{\vect{k}}b_{\vect{k}}$ this is $-\Delta_{\mathrm{BCS}}O_{\mathsf{s}_1}(\vect{0})$, the pair term of $H_{\mathrm{MF}}$.

\subsection*{Ward-identity violation of the bare bubble}

The quasiparticle bubble with vertices $A$ and $B$ follows from Wick's theorem for the Nambu field as in \cref{app:D.10}. It carries a trace over the Nambu index:
\begin{equation*}
  R^0_{AB}(\vect{q},\ii\omega^{\mathrm{B}}) = \frac1\Vol\sum_{\vect{k}}T\sum_{j}\Tr\bigl[A\,\GF_0(\vect{k}+\vect{q},\ii\omega^{\mathrm{F}}_j + \ii\omega^{\mathrm{B}})\,B\,\GF_0(\vect{k},\ii\omega^{\mathrm{F}}_j)\bigr] .
\end{equation*}
With the spectral form of $\GF_0$, the frequency sum is \cref{eq:D.10-sum} with $\xi_{\vect{k}} \to s'E_{\vect{k}}$ and $\xi_{\vect{k}+\vect{q}} \to sE_{\vect{k}+\vect{q}}$. At $T = 0$, $n_{\mathrm{F}}(-E) = 1$ and $n_{\mathrm{F}}(E) = 0$, and after continuation
\begin{equation}
  R^0_{AB}(\vect{q},\omega) = \frac1\Vol\sum_{\vect{k}}\left[\frac{\Tr[A P^+_{\vect{k}+\vect{q}}BP^-_{\vect{k}}]}{\omega - E_{\vect{k}} - E_{\vect{k}+\vect{q}} + \ii0^+} - \frac{\Tr[A P^-_{\vect{k}+\vect{q}}BP^+_{\vect{k}}]}{\omega + E_{\vect{k}} + E_{\vect{k}+\vect{q}} + \ii0^+}\right] .
  \label{eq:D.21-R0}
\end{equation}
For $A = B = \boldsymbol{\eta}$, $\boldsymbol{\eta}h_{\vect{k}'}\boldsymbol{\eta} = \xi'\boldsymbol{\eta} + \Delta_{\mathrm{BCS}}\mathsf{s}_1$, and the trace is $\frac12[1 - (\xi\xi' - \Delta_{\mathrm{BCS}}^2)/EE']$, the case-I creation factor of \cref{eq:D.20-cohE}, as it must be.

At $\vect{q} = 0$ the traces follow from the Pauli algebra of the Nambu matrices. Write $\mathsf{s}_3 \equiv \boldsymbol{\eta}$, so that $h_{\vect{k}}/E_{\vect{k}} = \sum_an_a\mathsf{s}_a$ with the unit vector $\vect{n} = (-\Delta_{\mathrm{BCS}}/E, 0, \xi/E)$. Multiplying out with $\mathsf{s}_a\mathsf{s}_b = \delta_{ab} + \ii\epsilon_{abc}\mathsf{s}_c$,
\begin{equation*}
  \Tr[\mathsf{s}_aP^+\mathsf{s}_bP^-] = \delta_{ab} - n_an_b - \ii\epsilon_{abc}n_c, \qquad \Tr[\mathsf{s}_aP^-\mathsf{s}_bP^+] = \delta_{ab} - n_an_b + \ii\epsilon_{abc}n_c .
\end{equation*}
With $F(\omega) = \Vol^{-1}\sum_{\vect{k}}[E_{\vect{k}}(\omega^2 - 4E_{\vect{k}}^2)]^{-1}$, the symmetric parts combine into $4E/(\omega^2 - 4E^2)$ and the antisymmetric parts into $2\omega/(\omega^2 - 4E^2)$, and
\begin{equation}
  \begin{gathered}
  R^0_{\boldsymbol{\eta}\boldsymbol{\eta}}(0,\omega) = 4\Delta_{\mathrm{BCS}}^2F, \qquad
  R^0_{\mathsf{s}_2\mathsf{s}_2}(0,\omega) = \omega^2F - \frac{2}{V_{\mathrm{BCS}}}, \\
  R^0_{\mathsf{s}_2\boldsymbol{\eta}}(0,\omega) = -R^0_{\boldsymbol{\eta}\mathsf{s}_2}(0,\omega) = 2\ii\Delta_{\mathrm{BCS}}\,\omega F .
  \end{gathered}
  \label{eq:D.21-q0}
\end{equation}
In $R^0_{\mathsf{s}_2\mathsf{s}_2}$ we wrote $4E/(\omega^2 - 4E^2) = [\omega^2 - (\omega^2 - 4E^2)]/[E(\omega^2 - 4E^2)]$ and used the gap equation at $T = 0$, $\Vol^{-1}\sum_{\vect{k}}1/E_{\vect{k}} = 2/V_{\mathrm{BCS}}$ (\cref{eq:D.20-gap}). The elements that involve $\mathsf{s}_1$ with $\boldsymbol{\eta}$ or $\mathsf{s}_2$ are odd in $\xi$ and vanish when the density of states is symmetric about the Fermi level over the shell, which we assume; the amplitude channel then decouples.

The bare density response at $\vect{q} = 0$ is therefore not zero. Its spectral weight sits at $\omega = 2E_{\vect{k}}$, with $S(0,\omega) = -2\Imag R^0 = (2\pi/\Vol)\sum_{\vect{k}}(\Delta_{\mathrm{BCS}}^2/E_{\vect{k}}^2)\,\delta(\omega - 2E_{\vect{k}})$, and its first moment is
\begin{equation}
  \int_0^\infty\dd\omega\,\omega\,S^0(0,\omega) = \frac{4\pi}{\Vol}\sum_{\vect{k}}\frac{\Delta_{\mathrm{BCS}}^2}{E_{\vect{k}}} = \frac{8\pi\Delta_{\mathrm{BCS}}^2}{V_{\mathrm{BCS}}} ,
  \label{eq:D.21-violation}
\end{equation}
by the gap equation once more. The f-sum rule \cref{eq:4.3-fsum} requires $\pi\nd_eq^2/m_e$, which vanishes at $q = 0$. The origin is visible in the commutator. With $[\mathsf{s}_1, \boldsymbol{\eta}] = -2\ii\mathsf{s}_2$,
\begin{equation*}
  [H_{\mathrm{MF}}, n_e(\vect{q})] = \vect{q}\cdot\vect{j}_e(\vect{q}) + 2\ii\Delta_{\mathrm{BCS}}\,O_{\mathsf{s}_2}(\vect{q}) ,
\end{equation*}
where the kinetic term gives $\sum_{\vect{k}}(\xi_{\vect{k}+\vect{q}} - \xi_{\vect{k}})(c^\dagger_{\vect{k}+\vect{q}\uparrow}c_{\vect{k}\uparrow} + c^\dagger_{\vect{k}+\vect{q}\downarrow}c_{\vect{k}\downarrow}) = \vect{q}\cdot\vect{j}_e(\vect{q})$ and the pair term the rest. The mean-field Hamiltonian converts density into pair phase, and charge leaks into the condensate without coming back.

\subsection*{The conserving response}

The response that conserves charge is the linear response of the time-dependent mean field. Let an external potential $\phi_{\mathrm{ext}}$ couple to $n_e(\vect{q})$. The quasiparticles then see the external field and the change of the pair field it induces, $\delta H = -(V_{\mathrm{BCS}}/2\Vol)\sum_a\delta\langle O_{\mathsf{s}_a}(-\vect{q})\rangle O_{\mathsf{s}_a}(\vect{q})$, which is the mean-field term of $H_{\mathrm{int}}$ at wavevector $\vect{q}$, and respond to both with the bare bubbles. In the channel basis $(\boldsymbol{\eta}, \mathsf{s}_1, \mathsf{s}_2)$ this is a random-phase approximation with the kernel $K = \operatorname{diag}(0, -V_{\mathrm{BCS}}/2, -V_{\mathrm{BCS}}/2)$:
\begin{equation}
  R = R^0 + R^0KR, \qquad
  R_{\boldsymbol{\eta}\boldsymbol{\eta}} = R^0_{\boldsymbol{\eta}\boldsymbol{\eta}} + \sum_{a,b=1,2}R^0_{\boldsymbol{\eta}\mathsf{s}_a}\Bigl[\bigl(1 - K R^0\bigr)^{-1}K\Bigr]_{ab}R^0_{\mathsf{s}_b\boldsymbol{\eta}} ,
  \label{eq:D.21-Rc}
\end{equation}
with all bubbles from \cref{eq:D.21-R0} and the inverse taken in the $(\mathsf{s}_1,\mathsf{s}_2)$ block. In diagrams, the correction term is the bubble with the ladder of pair propagators between its two vertices~\cite{anderson1958random}.

That this response obeys the continuity equation follows from a property of mean-field dynamics. Under $H_{\mathrm{MF}} + \delta H$ the state stays a quasi-free state, one for which Wick's theorem holds, at every time. For any one-body operator $A$, the average $\langle[H_0 + H_{\mathrm{int}}, A]\rangle$ in such a state, evaluated with Wick's theorem, equals $\langle[H_{\mathrm{MF}} + \delta H, A]\rangle$, provided the mean fields are the self-consistent ones: each pair of contractions of $H_{\mathrm{int}}$ is a mean field acting on the remaining two operators. The Hartree mean fields of $H_{\mathrm{int}}$, which are density--density terms, commute with $n_e$ and may be kept or dropped without changing the argument; we drop them, since their static part only shifts $\mu_{\mathrm{ch}}$. Hence $\dd\langle n_e(\vect{q})\rangle/\dd t = \ii\langle[H_0 + H_{\mathrm{int}}, n_e(\vect{q})]\rangle = \ii\vect{q}\cdot\langle\vect{j}_e(\vect{q})\rangle$, since $H_{\mathrm{int}}$ commutes with the density. The term $2\ii\Delta_{\mathrm{BCS}}\delta\langle O_{\mathsf{s}_2}\rangle$ that the bare dynamics produces is cancelled, to linear order, by the commutator of $n_e$ with $\delta H$, which contains $\langle O_{\mathsf{s}_1}(\vect{0})\rangle\delta\langle O_{\mathsf{s}_2}\rangle$ with the opposite sign. A response that obeys continuity at all times satisfies the f-sum rule, by the double-commutator argument of \cref{sec:4.3}, with the averages taken in the mean-field ground state, and it vanishes at $\vect{q} = 0$, where $n_e$ is the conserved particle number.

The explicit algebra at $\vect{q} = 0$ confirms the cancellation. With \cref{eq:D.21-q0}, the denominator is $1 - K_{22}R^0_{\mathsf{s}_2\mathsf{s}_2} = 1 + \frac12V_{\mathrm{BCS}}(\omega^2F - 2/V_{\mathrm{BCS}}) = \frac12V_{\mathrm{BCS}}\,\omega^2F$, and
\begin{equation*}
  R_{\boldsymbol{\eta}\boldsymbol{\eta}}(0,\omega) = 4\Delta_{\mathrm{BCS}}^2F + \frac{(-2\ii\Delta_{\mathrm{BCS}}\omega F)(-\frac12V_{\mathrm{BCS}})(2\ii\Delta_{\mathrm{BCS}}\omega F)}{\frac12V_{\mathrm{BCS}}\,\omega^2F} = 4\Delta_{\mathrm{BCS}}^2F - 4\Delta_{\mathrm{BCS}}^2F = 0 ,
\end{equation*}
for every $\omega$. The phase fluctuations of the condensate absorb exactly the spurious pair-breaking weight. The cancellation used the gap equation; a response built on a gap that does not solve it would not conserve charge.

\subsection*{The Goldstone mode}

The same denominator vanishes at $\omega = 0$ for $\vect{q} = 0$, because the gap equation makes $1 - K_{22}R^0_{\mathsf{s}_2\mathsf{s}_2}(0,0) = 0$: a uniform rotation of the phase of $\Delta_{\mathrm{BCS}}$ costs no energy. At small $\vect{q}$ the zero moves to a pole at $\omega = c_{\mathrm{AB}}q$, the Anderson--Bogoliubov mode, a sound wave of the condensate phase that carries the density response of a neutral superconductor at $\omega < 2\Delta_{\mathrm{BCS}}$ and exhausts the f-sum rule as $q \to 0$. Anderson found $c_{\mathrm{AB}} = v_{\mathrm{F}}/\sqrt3$~\cite{anderson1958random}. The zero of $\det(1 - KR^0)$, evaluated numerically from \cref{eq:D.21-R0} for a parabolic band at $q = 0.02\,\Delta_{\mathrm{BCS}}/v_{\mathrm{F}}$, lies at $c_{\mathrm{AB}} = 0.579\,v_{\mathrm{F}}$, against $1/\sqrt3 = 0.577$. With the long-range Coulomb interaction, which adds the random-phase kernel $K_{\boldsymbol{\eta}\boldsymbol{\eta}} = v_{\mathrm{C}}(q)$ in the $\boldsymbol{\eta}$ channel, the mode is pushed up to the plasma frequency, and the low-frequency density response of a charged superconductor is screened as in the normal metal~\cite{anderson1958random}. Adding that channel keeps the response conserving, since a density--density interaction commutes with the density.

\subsection*{The Ward identity}

Apply the continuity equation inside the three-point function $\langle\mathrm{T}\,\Psi_{\vect{k}+\vect{q}}\Psi^\dagger_{\vect{k}}\,n_e(\vect{q})^\dagger\rangle$, the amplitude for the density to scatter an electron from $\vect{k}$ to $\vect{k}+\vect{q}$. Its time derivative brings down the current, and the equal-time commutators of $n_e$ with the two fields give two terms with a single Green's function each. Multiply the whole by $\GF^{-1}$ on each side, which removes the propagation of the incoming and the outgoing electron and leaves the vertex alone (the step called amputation). The result is an identity for the dressed density vertex $\mathsf{V}_{\boldsymbol{\eta}}$ and current vertex $\boldsymbol{\mathsf{V}}_{\vect{j}}$:
\begin{equation}
  \vect{q}\cdot\boldsymbol{\mathsf{V}}_{\vect{j}}(\vect{k}+\vect{q},\vect{k}) - \omega\,\mathsf{V}_{\boldsymbol{\eta}}(\vect{k}+\vect{q},\vect{k}) = \boldsymbol{\eta}\,\GF^{-1}(\vect{k},\nu) - \GF^{-1}(\vect{k}+\vect{q},\nu + \omega)\,\boldsymbol{\eta} .
  \label{eq:D.21-ward}
\end{equation}
In the normal state, with $\GF^{-1} = \nu - \xi_{\vect{k}}\boldsymbol{\eta}$, the right side is $-\omega\boldsymbol{\eta} + (\xi_{\vect{k}+\vect{q}} - \xi_{\vect{k}})$, and the bare vertices $\mathsf{V}_{\boldsymbol{\eta}} = \boldsymbol{\eta}$ and $\boldsymbol{\mathsf{V}}_{\vect{j}} = (\vect{k} + \vect{q}/2)/m_e$ satisfy it for a parabolic band. In the superconductor, $\GF^{-1} = \nu - h_{\vect{k}}$ adds $\Delta_{\mathrm{BCS}}(\boldsymbol{\eta}\mathsf{s}_1 - \mathsf{s}_1\boldsymbol{\eta}) = 2\ii\Delta_{\mathrm{BCS}}\mathsf{s}_2$ to the right side, which no bare vertex supplies. The ladder supplies it. Its density vertex is $\mathsf{V}_{\boldsymbol{\eta}} = \boldsymbol{\eta} + \mathsf{c}_2\mathsf{s}_2$ with the phase-channel coefficient $\mathsf{c}_2 = K_{22}R^0_{\mathsf{s}_2\boldsymbol{\eta}}/(1 - K_{22}R^0_{\mathsf{s}_2\mathsf{s}_2})$, which at $\vect{q} = 0$ is, by \cref{eq:D.21-q0}, $(-\frac12V_{\mathrm{BCS}})(2\ii\Delta_{\mathrm{BCS}}\omega F)/(\frac12V_{\mathrm{BCS}}\omega^2F) = -2\ii\Delta_{\mathrm{BCS}}/\omega$. Then $-\omega\mathsf{V}_{\boldsymbol{\eta}} = -\omega\boldsymbol{\eta} + 2\ii\Delta_{\mathrm{BCS}}\mathsf{s}_2$, which is \cref{eq:D.21-ward} at $\vect{q} = 0$. The pole of $\mathsf{c}_2$ at $\omega = 0$ is the Goldstone mode; at finite $\vect{q}$ it moves to $\omega = c_{\mathrm{AB}}q$, and the vertex keeps satisfying the identity. The bare vertex is enough for the spin density $s^z_e$, whose Nambu vertex is the unit matrix: $\mathbb{1}\,h_{\vect{k}} - h_{\vect{k}}\,\mathbb{1} = 0$ in the gap term, so no vertex correction is needed, in agreement with \cref{app:D.20}.

A numerical check on a square-lattice model with this interaction, $48\times48$ sites and $\Delta_{\mathrm{BCS}} = 0.30t$, confirms the derivation. At $\vect{q} = 0$ the conserving $R_{\boldsymbol{\eta}\boldsymbol{\eta}}$ vanishes at frequencies below, inside and above the pair continuum, while the bare one does not. At the two smallest nonzero wavevectors of the lattice, the coefficient of $1/\omega^2$ in the conserving response, which is $\pi^{-1}$ times the first moment, reproduces the double commutator $\Vol^{-1}\sum_{\vect{k}\sigma}(\varepsilon_{\vect{k}+\vect{q}} + \varepsilon_{\vect{k}-\vect{q}} - 2\varepsilon_{\vect{k}})\langle n_{\vect{k}\sigma}\rangle$, the lattice form of $\nd_eq^2/m_e$. The bare bubble overestimates it by factors of 28 and 4, and the error grows as $q \to 0$ like the constant \cref{eq:D.21-violation} divided by $q^2$.
